\section{The ASML Question}\label{sec:background}

\subsection{A chokepoint}

Every leading-edge logic chip, including the accelerators on which frontier AI runs, is patterned with EUV lithography, as are the newest DRAM nodes, and ASML is the only firm that makes EUV machines. It recognized 48 EUV systems in 2025, four of them the newer High-NA version, and expected to ship about 65 low-NA systems in 2026. In October 2026 its chief executive said that planned supply for both 2027 and 2028 would not fully meet customer demand \citep{asml2026q3}. The machines cannot simply be bought and left alone. Installed-base management, ASML's term for servicing, upgrades, and spare parts, made up a quarter of its 2025 revenue, and EUV tools reportedly stop working quickly without ASML parts \citep{asml2026annual, bloomberg2024disable}.

China cannot buy EUV machines, but it is ASML's largest market for the older immersion DUV tools with which Chinese fabs make 7-nanometre chips by multiple patterning. Customers in China accounted for 26 percent of ASML's sales in 2023, 36 percent in 2024, and 29 percent in 2025 \citep{asml2026annual}. ASML expects China to fall to about 20 percent of its sales in 2026, although China's share of its system sales rose from 14 percent in the second quarter to 23 percent in the third \citep{asml2026q2call, asml2026q3}. Several hundred Dutch immersion tools are believed to be running in China, and they require maintenance roughly every six months \citep{aei2026lithography}. China's own efforts to escape the chokepoint are real but slow. A prototype EUV machine, built in a Shenzhen laboratory with the help of former ASML engineers, was generating EUV light but had not produced working chips as of December 2025; Beijing's goal is working chips by 2028, and people close to the project told Reuters that 2030 was more realistic \citep{reuters2025euv}.

\subsection{A nested chokepoint}

ASML is a Dutch company, but the machine it sells is not a Dutch product in any simple sense. Carl Zeiss SMT in Germany is the sole supplier of its optics, and ASML says that a prolonged loss of Zeiss supply would ``effectively'' end its business. The EUV light source is developed at Cymer in San Diego and driven by laser amplifiers from Trumpf. At the end of 2025 ASML had 8,242 employees in the United States, nearly a fifth of its workforce, and its manufacturing sites include Wilton, Connecticut, and San Diego \citep{asml2026annual}. This geography gives Washington jurisdiction over the firm. When the United States restricted support for advanced Chinese fabs by US persons in October 2022, ASML instructed its American staff to stop servicing Chinese customers within a week \citep{cnbc2022asml}.

The United States has turned that jurisdiction into a standing instrument. A December 2024 rule extended the foreign direct product rule to semiconductor manufacturing equipment and set a zero de minimis threshold for certain chipmaking tools, lithography tools among them, that contain US-origin chips and are bound for China. It excluded exports by firms in thirty-three listed countries, among them the Netherlands, Japan, and 25 of the 27 EU member states, through a supplement that the United States can amend at will \citep{bis2024fdpr}. In April 2026 a bipartisan bill, the MATCH Act, proposed a countrywide ban on sales of immersion lithography tools to China and restrictions on servicing at listed Chinese fabs, with 150 days for allies to show progress toward aligning before Washington would impose controls of its own \citep{baumgartner2026match}. As of late October 2026 it had not been enacted. In the terms of \citet{farrell2019weaponized}, ASML is a central node in an asymmetric network, but weaponizing a node requires effective jurisdiction over it, and here jurisdiction is shared: Washington's reach over ASML's US-origin subsystems, US staff, and US sites rivals that of The Hague.

\subsection{Three dependencies}

Europe's use of ASML is shaped by three dependencies, each of which the model treats explicitly.

\paragraph{Security.} Europe relies on the United States for nuclear deterrence, air and missile defence, intelligence, and enablers it cannot quickly replace. Bruegel and the Kiel Institute estimate that defending Europe without the United States would require about 300,000 additional troops and about EUR 250 billion a year in additional spending \citep{burilkov2025defending}. That sum is roughly eight times ASML's 2025 revenue and about 26 times its sales to China. The second Trump administration has tied economic rewards explicitly to burden-sharing and export-control alignment. Its National Security Strategy, released in December 2025, says the United States will stand ready to help, ``potentially through more favorable treatment on commercial matters, technology sharing, and defense procurement'', countries that take more responsibility for their own regional security and ``align their export controls with ours'' \citep{whitehouse2025nss}. It has also cast doubt on the guarantee itself. A drawdown of 25,000 or more US troops from Europe was under review in September 2026 \citep{nbc2026pentagon}, and in a poll published that June, majorities in all fifteen European countries surveyed doubted the United States would come to their defence \citep{ecfr2026home}.

\paragraph{Frontier AI.} Europe hosts about five percent of the world's AI compute and depends on American firms for frontier models \citep{juijn2026europe}. Access has already been treated as an instrument of policy. The January 2025 AI Diffusion Rule divided the world into tiers and placed ten EU members in the top tier and the other seventeen below it \citep{bis2025diffusion}. The August 2025 joint statement on the EU--US trade framework recorded the EU's intention to buy at least \$40 billion of US AI chips and to adopt technology-security requirements in line with those of the United States, and said Washington would facilitate those exports once the requirements were in place \citep{whitehouse2025jointstatement}.

\paragraph{China.} Europe imports almost all of its rare-earth permanent magnets from China, and Beijing has shown that it can use licensing to cut supply in weeks. Its October 2025 controls, which reached foreign products containing small amounts of Chinese rare earths, were suspended after the US--China meeting in Busan, not withdrawn \citep{mofcom2025suspension}. The Nexperia dispute is the closest real analogue to the scenario's endgame. On September 30, 2025, the Dutch government used emergency powers to take control of the Chinese-owned chipmaker after a US rule threatened its operations. China blocked exports from Nexperia's Chinese plants, which threatened production at European carmakers, and on November 19 the Dutch minister suspended the order \citep{taylorwessing2025nexperia}. Relief came mainly through the US--China deal, not through anything Europe obtained for itself.

\subsection{What the scenario claims}

\emph{Europe 2031} makes five claims that a model can test. (1) In 2028, faced with the threat that Washington will apply its foreign direct product rule to ASML, the Netherlands has no choice but to halt ASML's remaining exports and servicing of immersion DUV tools in China, and gets nothing for it. (2) In 2029 Europe cannot use ASML against the United States because security-dependent members and the Netherlands itself block the countermeasure. (3) In 2031 Washington moves to take control of ASML because it fears losing the chokepoint to a Europe drifting toward China. (4) Its threat of permanent exclusion from US AI, backed by a stockpile of EUV tools, lets it outlast Europe. (5) Compute on European soil and a coalition of middle powers would have given Europe leverage. The model confirms (1) and (3), rejects the mechanism behind (2), finds the threat in (4) incredible as stated, and specifies when (5) would and would not work.
